% Supplementary material of the survfcbf Application Note, compiled by supplementary.tex
% into the single file that goes both to the journal and, as a separate file, to bioRxiv.
% Every number in a table, and every figure quoted below through a \Supp... macro, is
% generated from the tracked result files by paper/make_supplementary_tables.R. Settings
% stated in words (folds, grids, seeds) are those of the scripts named in each section.

\noindent This supplement documents the analyses behind every result of the Application
Note. The scripts it names are in the folder \code{experiments/benchmark/realdata/} of
the repository, except \code{grid\_lib.R} and \code{sim\_matched\_size.R}, which are in
\code{experiments/benchmark/}.
Every table is generated from the result files tracked in the repository by
\code{paper/make\_supplementary\_tables.R}.

\section{Machine-learning reporting}\label{S-sec:ml}

The characterisation in the Application Note trains prognostic (risk) models, so this section
gives the description that \emph{Bioinformatics} asks of machine-learning papers: the
composition of the training data, the preparation of the cross-validation sets and the
performance on test data that no training step has seen.

\paragraph{What is trained.} \pkg{survfcbf} is a feature selector, and the package ships
no fitted prognostic model. The object under evaluation is a procedure: a selector
applied to a training set, followed by a ridge-penalised Cox model \citep{simon2011}
fitted on the selected features of the same training set. Every component of the
procedure is refitted from scratch on each training set, and no fitted model is reused
across folds or across cohorts. The models are prognostic: they rank patients by their risk
of death under the treatments they received, and nothing here addresses the prediction of
benefit from a treatment.

\paragraph{Training data.} The data are the eleven TCGA cohorts of the reference
benchmark \citep{bommert2021bib}, with \SuppTcgaPatients{} patients and
\SuppTcgaEvents{} events in total, and METABRIC, with \SuppMetabricN{} patients and
\SuppMetabricEvents{} events (Table~\ref{S-tab:cohorts}). For TCGA the features are all
the RNA and miRNA expression features of the OpenML datasets prepared by
\citet{herrmann2020}, the feature set of the reference benchmark, and the outcome is
overall survival: \citet{herrmann2020} define it as the censored survival time with
deaths as events. The features are used as distributed, without a further filter or
transformation. OpenML lists these datasets under GPL-2 and does not document their
normalisation. For METABRIC the features are the genes of the complete Illumina HT-12 v3
profile, less those with a missing value or zero variance, and the outcome is overall
survival. The 19{,}850 profiled genes and 1{,}980 samples become \SuppMetabricP{} genes
and the \SuppMetabricN{} patients with a clinical record and a positive overall survival
time. The gene filters do not use the outcome, and no other preprocessing precedes the
cross-validation. The loaders check every input file against recorded checksums (SHA-256
for the METABRIC snapshot, MD5 for the OpenML files) and stop on any mismatch.

\paragraph{Batch effects.} No batch correction is applied. Every comparison is between
selectors within one cohort, on the same folds and the same expression matrix, the input
of the reference benchmark. A technical effect is therefore shared by all selectors, but it
can still favour some of them, because it shapes what each criterion picks. A ranking by
variance, which ignores the outcome, follows the largest sources of variation, technical
ones included, and the same holds for its stability and for the enrichment of its ranking
(Section~\ref{S-sec:gsea}). The tissue source site, read from the TCGA barcodes
(\SuppCentreMin{} to \SuppCentreMax{} centres per cohort), predicts censoring with a
cross-validated Harrell's C of \SuppCentreCens{} (\SuppCentreCensMin{} to
\SuppCentreCensMax{} across cohorts) and death with \SuppCentreDeath, and the follow-up of
the censored patients differs between centres in \SuppCentreDiffer{} of eleven cohorts
(Kruskal-Wallis $P < 0.001$, \code{center\_followup.R}). Follow-up therefore varies with
something that may also shape expression, and Section~\ref{S-sec:selectors} tests whether
the selected panels inherit it. METABRIC, on a microarray platform, is analysed on its
own, so no correction between platforms is needed.

\begin{table}[htbp]
\centering\footnotesize
\caption{Cohorts. $n$ is the number of patients analysed and $p$ the number of features
in the main benchmark. ``$p$ coding'' counts the protein-coding RNA features, the
feature set of the matched panel size analysis (Section~\ref{S-sec:cardinality}) and of
the gene-set enrichment (Section~\ref{S-sec:gsea}). Each TCGA cohort is assembled from
the OpenML datasets whose identifiers are listed as its data source, as in
\citet{herrmann2020} and \citet{bommert2021bib}.
METABRIC is the profile \code{brca\_metabric\_mrna} of the cBioPortal study
\code{brca\_metabric}, archived at \href{https://doi.org/\metabricdoi}{doi:\metabricdoi}.
n.r., not run.}
\label{S-tab:cohorts}
\begin{tabularx}{\linewidth}{@{}l >{\raggedright\arraybackslash}X l rrrrr@{}}
\toprule
Cohort & Cancer type & Data source & $n$ & Events & Censored (\%) & $p$ & $p$ coding \\
\midrule
BLCA & Bladder urothelial carcinoma & 42279, 42280 & 382 & 103 & 73.0 & 23{,}906 & 16{,}390 \\
BRCA & Breast invasive carcinoma & 42323, 42324, 42325 & 735 & 72 & 90.2 & 23{,}529 & 16{,}307 \\
HNSC & Head and neck squamous cell carcinoma & 42285, 42286 & 443 & 152 & 65.7 & 22{,}313 & 16{,}202 \\
KIRC & Kidney renal clear cell carcinoma & 42287, 42288 & 249 & 62 & 75.1 & 23{,}697 & 16{,}206 \\
LGG & Brain lower grade glioma & 42293, 42294 & 419 & 77 & 81.6 & 22{,}942 & 15{,}993 \\
LUAD & Lung adenocarcinoma & 42297, 42298 & 426 & 101 & 76.3 & 24{,}480 & 16{,}527 \\
LUSC & Lung squamous cell carcinoma & 42299, 42300 & 418 & 132 & 68.4 & 24{,}419 & 16{,}595 \\
OV & Ovarian serous cystadenocarcinoma & 42301, 42302 & 219 & 109 & 50.2 & 25{,}483 & 16{,}492 \\
PAAD & Pancreatic adenocarcinoma & 42303, 42304 & 124 & 52 & 58.1 & 22{,}960 & 16{,}386 \\
SKCM & Skin cutaneous melanoma & 42311, 42312 & 249 & 87 & 65.1 & 23{,}250 & 16{,}115 \\
STAD & Stomach adenocarcinoma & 42313, 42314 & 295 & 62 & 79.0 & 26{,}814 & 16{,}899 \\
\midrule
TCGA & Eleven cohorts &  & 3{,}959 & 1{,}009 & 74.5 &  &  \\
METABRIC & Breast cancer & cBioPortal & 1{,}979 & 1{,}143 & 42.2 & 19{,}839 & n.r. \\
\bottomrule
\end{tabularx}

\end{table}

\paragraph{Cross-validation sets.} Each cohort is split into five outer folds, stratified
on the event indicator. Within the patients with an event, and within the censored
patients, fold sizes differ by at most one, so every test fold holds a fifth of the
patients and a fifth of the events up to rounding. The assignment depends only on the
event indicator and a fixed seed, so within a cohort every selector is evaluated on the
same five splits. In each fold the other four fifths form the training set, and every
quantity learned from data is learned on it alone (Table~\ref{S-tab:leakage}). The
penalties of the lasso and of the ridge Cox models are tuned by an inner cross-validation
inside the training set (the ten-fold internal cross-validation of \code{cv.glmnet}), and
so is $q$ in the tuned variant of \pkg{survfcbf} of the simulation (three inner folds,
Section~\ref{S-sec:simulation}). Leave-one-out cross-validation is not used anywhere.

\begin{table}[htbp]
\centering\footnotesize
\caption{Where each quantity learned from data is fitted. Training set: the four fifths
of the cohort outside the outer test fold. Inner cross-validation: folds drawn inside
that training set. The test fold enters only the predictions being scored.}
\label{S-tab:leakage}
\begin{tabularx}{\linewidth}{@{}X l l@{}}
\toprule
Quantity & Used by & Fitted on \\
\midrule
Null Cox model and its martingale residuals & \pkg{survfcbf}, MIM, mRMR & Training set \\
Equal-frequency discretisation cut-points & \pkg{survfcbf}, MIM, mRMR & Training set \\
Relevance threshold $\delta$ from the permutation null & \pkg{survfcbf} & Training set \\
Number of bins $q$ & \pkg{survfcbf}, tuned $q$ (simulation) & Inner cross-validation \\
Feature variances & Variance & Training set \\
CAR scores and their censoring weights & CAR & Training set \\
Lasso penalty and coefficients & Lasso Cox & Inner cross-validation \\
Ridge penalty and coefficients of the downstream model & All selectors & Inner cross-validation \\
Censoring distribution for Uno's C and the integrated Brier score & Metrics & Training set \\
Truncation time of Uno's C (80th percentile of follow-up) & Metrics & Training set \\
Feature correlations for the stability index & Metrics & All patients, without the outcome \\
Feature correlations for the redundancy within a panel & Metrics & Training set \\
\bottomrule
\end{tabularx}
\end{table}

The one quantity computed on all patients is the feature correlation matrix of the
stability index. It does not use the outcome and is only used to score the agreement
between sets that were already selected, so it cannot influence any selection or
prediction.

\paragraph{Independent test sets.} In each fold the held-out fifth is the independent
test set. None of its patients enters selection, tuning, model fitting or the censoring
weights of that fold. Uno's C and the integrated Brier score are reported for every
cohort as the mean and standard error over the five test folds of the published split
(Tables~\ref{S-tab:unoc} and~\ref{S-tab:ibs}), not only as an average over cohorts. These
standard errors treat the five folds as independent, which they are not, since their
training sets overlap, so they understate the uncertainty of a single cohort's estimate.
The comparisons quoted in the Application Note average ten such splits per cohort
instead (Section~\ref{S-sec:bayes}). METABRIC is an
independent breast cancer cohort measured
on a different platform. It is analysed with the same protocol rather than used as a
test set for models trained on TCGA, because what the characterisation compares is a
selection procedure, and no fitted model is transferred between cohorts.

\paragraph{Failed folds and null fits.} A fold fails when the selector returns fewer
features than the learner needs (the ridge Cox model of \pkg{glmnet} needs at least two)
or when the learner stops with an error. A failed fold has no prediction, is not scored
and is counted (Table~\ref{S-tab:failures}). A fold is a null fit when the fitted model
predicts the same survival curve for every test patient, which happens when the inner
cross-validation of the ridge model selects its largest penalty. Null fits are scored and
counted. The ridge penalty is taken at the minimum of the inner cross-validation error
(\code{lambda.min}) in every run, because the one-standard-error rule returns the null
model whenever events are scarce.

\paragraph{Seeds.} Every run is reproducible from one master seed. In the real-data
analyses the seed is 1 for every cohort, and the seed of each fold is derived from it
arithmetically (option \code{rng = "sequential"}, the package default). In the
simulation, replicate $r$ of every configuration is generated in its own
L'Ecuyer-CMRG stream seeded with $1000 + r$, and its folds are seeded from the same
number. The package also offers an independent L'Ecuyer-CMRG stream per fold
(\code{rng = "stream"}), which these runs did not use. The repeated cross-validation uses
seeds 1 to 10, one per repetition, seed 1 being the published split.

\section{Selectors and metrics}\label{S-sec:selectors}

Table~\ref{S-tab:selectors} lists the selectors and their settings, as configured in
\code{run\_dataset.R}. Every selector is followed by the
same downstream model, a ridge Cox model fitted with \code{cv.glmnet} through
\pkg{mlr3extralearners} (\code{alpha = 0}, \code{s = "lambda.min"}). The rankers keep the
top ten per cent of the features, the conventional cut-off.

\begin{table}[htbp]
\centering\footnotesize
\caption{Selectors. SU, symmetrical uncertainty. The martingale residuals come from a
covariate-free Cox model fitted on the training set. The rule of \citet{bommert2021bib}
sets $q = \max(\min(\lfloor n/3 \rfloor, 10), 2)$, which gives $q = 10$ at every training
set size in this study.}
\label{S-tab:selectors}
\begin{tabularx}{\linewidth}{@{}l X l@{}}
\toprule
Selector & Criterion & Panel size \\
\midrule
\pkg{survfcbf}, top-$K$ & SU with the martingale residuals at $q = 2$, ties broken by the
  absolute Spearman correlation of the continuous feature with the residuals, then
  Markov-blanket pruning until $K$ predominant features remain
  (\code{survfcbf\_select(q = 2, top = 10)}) & $K = 10$ \\
\pkg{survfcbf}, permutation & SU at $q$ from the rule of \citeauthor{bommert2021bib}, relevance threshold
  $\delta$ at the 95th percentile of the maximum SU over 200 permutations of the target, the
  same tie-break, then Markov-blanket pruning & Self-selected \\
Variance & Variance of each feature, without the outcome & 10\% of $p$ \\
CAR & Correlation-adjusted regression survival scores \citep{welchowski2019}, \pkg{carSurv} & 10\% of $p$ \\
MIM & Mutual information with the discretised martingale residuals \citep{kursa2021}, \pkg{praznik} & 10\% of $p$ \\
Lasso Cox & Non-zero coefficients of an $L_1$-penalised Cox model at \code{lambda.min}
  \citep{simon2011} & Self-selected \\
\bottomrule
\end{tabularx}
\end{table}

\paragraph{Metrics.} Discrimination is Uno's C \citep{uno2011}, the concordance index
with inverse probability of censoring weights (\code{surv.cindex} with
\code{weight\_meth = "G2"} in \pkg{mlr3proba}, \citealp{sonabend2021}), truncated at
the 80th percentile of the follow-up times of the training set (\code{t\_max}, computed
by \code{uno\_truncated.R}). Without truncation its censoring weights grow without bound
for test events beyond the last censoring time of the training set: over the
\SuppTruncFolds{} folds of \code{panel\_redundancy.R}, the untruncated value differed from
the truncated one by more than 0.10 in \SuppTruncOffTen\% of folds and by more than 0.20
in \SuppTruncOffTwenty\%. The untruncated value is still reported. Discrimination and
calibration together are the integrated Brier score of \citet{graf1999}
(\code{surv.graf}). Both take the censoring distribution from the training set of the
fold. Stability is the similarity-adjusted intersection count of
\pkg{stabm} \citep{bommert2021stabm} over the sets selected in the outer folds of a
cohort, counting only folds that selected at least one feature, with two
features counted as similar when the absolute Pearson correlation of their expression is
at least 0.9, so that exchanging a gene for a near-copy is not counted as instability.
Its correction for chance is estimated by simulation, the \pkg{stabm} default, and it
needs at least two non-empty sets. Panel size is the mean number of features selected per
fold, over the folds that selected at least one.

\paragraph{Ties in relevance.} At $q = 2$ a feature and the target are both split at
their medians, so SU depends on a single cell of a two-by-two table and takes few
distinct values: between \SuppDistinctSuMin{} and \SuppDistinctSuMax{} over the
\SuppCodingMin{} to \SuppCodingMax{} protein-coding genes of a cohort. Many genes therefore
tie on relevance, and the order in which tied genes are examined decides which of them
enter the panel. \pkg{survfcbf} breaks these ties by the absolute Spearman correlation
between the continuous feature and the martingale residuals, computed on the training
set, which makes the selection independent of the column order of the data. A test in
the package checks this invariance. With the earlier rule, which broke ties by column
position, permuting the columns of the same simulated data changed the top-$K$ panel.

\paragraph{What the $q = 2$ target encodes.} The martingale residual of patient $i$ is
$\delta_i - \hat H_0(t_i)$. While the cumulative baseline hazard stays below one, every
event has a positive residual and every censored patient one at or below zero, and among the
censored the residual is highest for the shortest follow-up. Splitting at the median
therefore puts in the upper class most of the events (all of them in cohorts where the
hazard never accumulates to one) and, when events are fewer than half the patients, the
censored patients followed for the shortest time (Table~\ref{S-tab:target}). The relevance of a gene then reflects follow-up as well as
risk. In TCGA, follow-up depends on the centre and the date of inclusion, which may go
with technical variation in expression, so this is one way a technical effect could
enter the relevance.

\paragraph{Whether the panel inherits follow-up.} \code{panel\_redundancy.R} tests whether
this reaches the selected genes. Each ten-gene panel is chosen on the true outcome, and the
downstream model is then trained and scored with the status reversed, so that censoring is
the event. Its Harrell's C measures how well the panel predicts the length of follow-up,
against 100 panels of ten random genes per cohort, each drawn from its own seeded stream
(Table~\ref{S-tab:censoring}). The top-$K$ panel at $q = 2$ scores \SuppCensTopk{} and
never exceeds the 95th percentile of the random panels (its highest percentile is
\SuppCensTopkMaxPct, in \SuppCensTopkMaxCohort), while top-$K$ at $q = 4$ and at $q = 10$
exceeds it in \SuppCensAboveQfour{} of eleven cohorts, MIM in \SuppCensAboveMim{} and mRMR
in \SuppCensAboveMrmr. In METABRIC (\code{metabric\_censoring.R}) the $q = 2$ panel sits at
the \SuppCensMetabricTopk{} percentile, and the $q = 10$ panel and mRMR at the
\SuppCensMetabricQten. The random panels themselves average \SuppCensRandomMean, so
follow-up is weakly associated with expression as a whole, in line with the centre effect
of Section~\ref{S-sec:ml}.

\paragraph{An alternative target.} \code{survfcbf\_select(target\_method =
"ipcw\_binned")} bins the observed time and weights patients by the inverse probability of
censoring, so censored patients carry no weight and follow-up does not enter the target
as above (\code{target\_sensitivity.R}). On the same folds it gives the same
discrimination (truncated Uno's C \SuppIpcwC{} against \SuppMartC), entirely different
panels (Jaccard index \SuppIpcwJaccardMax{} in every fold) and panels that predict censoring
at the level of random genes (\SuppIpcwCens, against \SuppMartCens{} for the martingale
target). Across cohorts, the censoring C of the martingale panels rises with the share of
censored patients in the upper class (Spearman \SuppGradMart), and that of the IPCW panels
less so (\SuppGradIpcw). It does not change BRCA (\SuppBrcaMartCens{} and
\SuppBrcaIpcwCens).

\begin{table}[htbp]
\centering\footnotesize
\caption{Cox model on censoring. Each ten-gene panel is selected on the true outcome over
the protein-coding features, and the ridge Cox model is trained and scored with censoring as
the event, on the folds of the matched panel size analysis. Harrell's C over the eleven
TCGA cohorts: mean, minimum and maximum. Percentile: mean over cohorts of the panel's
percentile among the 100 random ten-gene panels of its cohort. Above 95th: cohorts in which
the panel exceeds the 95th percentile. METABRIC: the same percentile in METABRIC. Null fits:
folds, out of 55 per panel (5,500 for the random panels), in which the model fell back to
the null model. n.r., not run.}
\label{S-tab:censoring}
\begin{tabular}{@{}l rrr rr r r@{}}
\toprule
 & \multicolumn{5}{c}{TCGA, eleven cohorts} &  &  \\
\cmidrule(lr){2-6}
Panel of ten genes & Mean & Min & Max & Percentile & Above 95th & METABRIC & Null fits \\
\midrule
\pkg{survfcbf}, $q = 2$ & 0.526 & 0.470 & 0.603 & 0.26 & 0 & 0.52 & 9 \\
\pkg{survfcbf}, $q = 4$ & 0.573 & 0.484 & 0.682 & 0.75 & 5 & n.r. & 1 \\
\pkg{survfcbf}, $q = 10$ & 0.576 & 0.459 & 0.667 & 0.84 & 5 & 1.00 & 0 \\
mRMR & 0.568 & 0.513 & 0.668 & 0.76 & 3 & 1.00 & 0 \\
mRMR, $q = 2$ & 0.540 & 0.482 & 0.615 & 0.43 & 1 & n.r. & 3 \\
MIM & 0.574 & 0.468 & 0.676 & 0.77 & 4 & 0.90 & 0 \\
MIM, $q = 2$ & 0.548 & 0.487 & 0.644 & 0.56 & 1 & n.r. & 5 \\
Variance & 0.509 & 0.457 & 0.658 & 0.16 & 1 & n.r. & 19 \\
CAR & 0.542 & 0.492 & 0.605 & 0.42 & 2 & n.r. & 1 \\
\midrule
100 random panels & 0.546 & 0.502 & 0.601 &  &  &  & 470 \\
\bottomrule
\end{tabular}

\end{table}

\begin{table}[htbp]
\centering\footnotesize
\caption{Composition of the upper class of the $q = 2$ target, computed on each whole
TCGA cohort. Censored: censored patients in the upper class, as a count and as a share of
the class. Follow-up, in days: median follow-up of the censored patients in the upper and
in the lower class, and median follow-up of the cohort by reverse Kaplan-Meier
(\SuppFollowMin{} to \SuppFollowMax{} days).}
\label{S-tab:target}
\begin{tabular}{@{}l rrr rr rrr@{}}
\toprule
 &  &  & \multicolumn{3}{c}{Upper class} & \multicolumn{3}{c}{Follow-up (days)} \\
\cmidrule(lr){4-6}\cmidrule(l){7-9}
Cohort & $n$ & Events & Size & Censored & Share (\%) & Upper & Lower & Cohort \\
\midrule
BLCA & 382 & 103 & 190 & 87 & 45.8 & 60 & 445 & 399 \\
BRCA & 735 & 72 & 365 & 293 & 80.3 & 168 & 1{,}150 & 554 \\
HNSC & 443 & 152 & 221 & 79 & 35.7 & 113 & 596 & 646 \\
KIRC & 249 & 62 & 120 & 60 & 50.0 & 123 & 1{,}433 & 1{,}165 \\
LGG & 419 & 77 & 205 & 133 & 64.9 & 117 & 650 & 475 \\
LUAD & 426 & 101 & 213 & 121 & 56.8 & 29 & 298 & 255 \\
LUSC & 418 & 132 & 209 & 88 & 42.1 & 25 & 505 & 578 \\
OV & 219 & 109 & 105 & 16 & 15.2 & 96 & 620 & 1{,}367 \\
PAAD & 124 & 52 & 62 & 14 & 22.6 & 12 & 330 & 488 \\
SKCM & 249 & 87 & 124 & 43 & 34.7 & 14 & 1{,}221 & 1{,}463 \\
STAD & 295 & 62 & 134 & 74 & 55.2 & 17 & 238 & 200 \\
\bottomrule
\end{tabular}

\end{table}

\section{Results per cohort}\label{S-sec:percohort}

Tables~\ref{S-tab:unoc} to \ref{S-tab:failures} give the result of every selector in every
cohort on the published split (seed 1), with the untruncated Uno's C. The panel sizes, failures and stability quoted in the Application Note
come from these tables. Its discrimination figures come from the ten-split analysis with
the truncated C (Section~\ref{S-sec:bayes}, Table~\ref{S-tab:rcvcohort}).

\begin{table}[htbp]
\centering\footnotesize
\caption{Uno's C per cohort, untruncated, on the published split: mean (standard error)
over the five outer test folds. An
asterisk marks results scored on fewer than five folds, because the other folds failed
(Table~\ref{S-tab:failures}), and a value without a standard error comes from a single
fold. n.m., no fold produced a model. n.r., not run. The means are over the cohorts with
a value, counted in the row below them.}
\label{S-tab:unoc}
\setlength{\tabcolsep}{4pt}
\begin{tabular}{@{}l rr rrrr@{}}
\toprule
 & \multicolumn{2}{c}{\pkg{survfcbf}} & \multicolumn{4}{c}{Baselines} \\
\cmidrule(lr){2-3}\cmidrule(l){4-7}
Cohort & top-$K$ & permutation & Variance & CAR & MIM & Lasso Cox \\
\midrule
BLCA & 0.571 (0.022) & n.m. & 0.649 (0.033) & 0.669 (0.026) & 0.669 (0.026) & 0.651 (0.027) \\
BRCA & 0.485 (0.041) & n.m. & 0.632 (0.046) & 0.623 (0.055) & 0.604 (0.040) & 0.591 (0.066)$^{\ast}$ \\
HNSC & 0.595 (0.026) & n.m. & 0.602 (0.031) & 0.618 (0.033) & 0.620 (0.040) & 0.618 (0.025) \\
KIRC & 0.521 (0.056) & 0.480$^{\ast}$ & 0.594 (0.073) & 0.648 (0.037) & 0.513 (0.118) & 0.614 (0.067) \\
LGG & 0.688 (0.056) & 0.730 (0.041)$^{\ast}$ & 0.746 (0.053) & 0.716 (0.032) & 0.726 (0.018) & 0.733 (0.030) \\
LUAD & 0.523 (0.025) & n.m. & 0.573 (0.036) & 0.569 (0.034) & 0.552 (0.040) & 0.542 (0.037) \\
LUSC & 0.486 (0.037) & n.m. & 0.508 (0.018) & 0.505 (0.013) & 0.513 (0.011) & 0.491 (0.011)$^{\ast}$ \\
OV & 0.478 (0.036) & n.m. & 0.525 (0.042) & 0.535 (0.039) & 0.540 (0.038) & 0.591 (0.072)$^{\ast}$ \\
PAAD & 0.599 (0.037) & 0.612 (0.011)$^{\ast}$ & 0.671 (0.023) & 0.650 (0.024) & 0.656 (0.010) & 0.616 (0.020) \\
SKCM & 0.511 (0.082) & 0.580$^{\ast}$ & 0.583 (0.074) & 0.612 (0.064) & 0.636 (0.061) & 0.642 (0.103)$^{\ast}$ \\
STAD & 0.511 (0.051) & 0.518 (0.043) & 0.593 (0.043) & 0.590 (0.044) & 0.583 (0.046) & n.m. \\
\midrule
Mean & 0.543 & 0.584 & 0.607 & 0.612 & 0.601 & 0.609 \\
Cohorts & 11 & 5 & 11 & 11 & 11 & 10 \\
\midrule
METABRIC & 0.573 (0.010) & n.r. & 0.614 (0.013) & 0.618 (0.006) & 0.615 (0.009) & 0.608 (0.006) \\
\bottomrule
\end{tabular}

\end{table}

\begin{table}[htbp]
\centering\footnotesize
\caption{Integrated Brier score per cohort (lower is better), mean (standard error) over
the five outer test folds. Symbols as in Table~\ref{S-tab:unoc}.}
\label{S-tab:ibs}
\setlength{\tabcolsep}{4pt}
\begin{tabular}{@{}l rr rrrr@{}}
\toprule
 & \multicolumn{2}{c}{\pkg{survfcbf}} & \multicolumn{4}{c}{Baselines} \\
\cmidrule(lr){2-3}\cmidrule(l){4-7}
Cohort & top-$K$ & permutation & Variance & CAR & MIM & Lasso Cox \\
\midrule
BLCA & 0.245 (0.026) & n.m. & 0.228 (0.043) & 0.229 (0.050) & 0.233 (0.056) & 0.234 (0.037) \\
BRCA & 0.177 (0.022) & n.m. & 0.171 (0.031) & 0.171 (0.031) & 0.170 (0.033) & 0.166 (0.036)$^{\ast}$ \\
HNSC & 0.168 (0.021) & n.m. & 0.165 (0.015) & 0.161 (0.014) & 0.166 (0.016) & 0.173 (0.018) \\
KIRC & 0.194 (0.015) & 0.156$^{\ast}$ & 0.180 (0.017) & 0.179 (0.016) & 0.181 (0.017) & 0.179 (0.015) \\
LGG & 0.161 (0.028) & 0.164 (0.016)$^{\ast}$ & 0.149 (0.021) & 0.149 (0.016) & 0.147 (0.014) & 0.141 (0.023) \\
LUAD & 0.158 (0.019) & n.m. & 0.150 (0.017) & 0.151 (0.017) & 0.159 (0.022) & 0.170 (0.022) \\
LUSC & 0.224 (0.013) & n.m. & 0.200 (0.006) & 0.202 (0.006) & 0.201 (0.006) & 0.219 (0.011)$^{\ast}$ \\
OV & 0.179 (0.019) & n.m. & 0.167 (0.013) & 0.172 (0.020) & 0.168 (0.016) & 0.150 (0.012)$^{\ast}$ \\
PAAD & 0.157 (0.038) & 0.115 (0.001)$^{\ast}$ & 0.121 (0.022) & 0.128 (0.026) & 0.130 (0.027) & 0.143 (0.024) \\
SKCM & 0.223 (0.035) & 0.168$^{\ast}$ & 0.202 (0.034) & 0.203 (0.037) & 0.198 (0.037) & 0.249 (0.053)$^{\ast}$ \\
STAD & 0.207 (0.024) & 0.202 (0.022) & 0.197 (0.025) & 0.200 (0.028) & 0.204 (0.026) & n.m. \\
\midrule
Mean & 0.190 & 0.161 & 0.175 & 0.177 & 0.178 & 0.183 \\
Cohorts & 11 & 5 & 11 & 11 & 11 & 10 \\
\midrule
METABRIC & 0.175 (0.005) & n.r. & 0.162 (0.005) & 0.163 (0.005) & 0.163 (0.004) & 0.172 (0.004) \\
\bottomrule
\end{tabular}

\end{table}

\begin{table}[htbp]
\centering\footnotesize
\caption{Selection stability per cohort, over the sets selected on the outer training
sets, counting only folds that selected at least one feature. n.m., fewer than two folds
selected a feature.}
\label{S-tab:stability}
\setlength{\tabcolsep}{4pt}
\begin{tabular}{@{}l rr rrrr@{}}
\toprule
 & \multicolumn{2}{c}{\pkg{survfcbf}} & \multicolumn{4}{c}{Baselines} \\
\cmidrule(lr){2-3}\cmidrule(l){4-7}
Cohort & top-$K$ & permutation & Variance & CAR & MIM & Lasso Cox \\
\midrule
BLCA & 0.090 & 0.000 & 0.944 & 0.422 & 0.180 & 0.264 \\
BRCA & 0.100 & 0.000 & 0.963 & 0.408 & 0.219 & 0.201 \\
HNSC & 0.100 & n.m. & 0.951 & 0.412 & 0.187 & 0.276 \\
KIRC & 0.190 & 0.000 & 0.944 & 0.438 & 0.165 & 0.214 \\
LGG & 0.119 & 0.000 & 0.954 & 0.442 & 0.303 & 0.243 \\
LUAD & 0.110 & n.m. & 0.946 & 0.404 & 0.188 & 0.193 \\
LUSC & 0.140 & 0.000 & 0.952 & 0.398 & 0.184 & 0.084 \\
OV & 0.090 & n.m. & 0.928 & 0.378 & 0.122 & 0.133 \\
PAAD & 0.050 & 0.000 & 0.914 & 0.392 & 0.160 & 0.154 \\
SKCM & 0.059 & 0.000 & 0.930 & 0.403 & 0.194 & 0.036 \\
STAD & 0.100 & 0.081 & 0.945 & 0.416 & 0.318 & 1.000 \\
\midrule
Mean & 0.104 & 0.010 & 0.943 & 0.410 & 0.202 & 0.254 \\
Cohorts & 11 & 8 & 11 & 11 & 11 & 11 \\
\midrule
METABRIC & 0.099 & n.r. & 0.979 & 0.409 & 0.592 & 0.359 \\
\bottomrule
\end{tabular}

\end{table}

\begin{table}[htbp]
\centering\footnotesize
\caption{Panel size per cohort, the mean number of features selected per outer fold,
over the folds that selected at least one. n.m., no fold selected a feature.}
\label{S-tab:panel}
\setlength{\tabcolsep}{4pt}
\begin{tabular}{@{}l rr rrrr@{}}
\toprule
 & \multicolumn{2}{c}{\pkg{survfcbf}} & \multicolumn{4}{c}{Baselines} \\
\cmidrule(lr){2-3}\cmidrule(l){4-7}
Cohort & top-$K$ & permutation & Variance & CAR & MIM & Lasso Cox \\
\midrule
BLCA & 10.0 & 1.0 & 2{,}391.0 & 2{,}391.0 & 2{,}391.0 & 17.6 \\
BRCA & 10.0 & 1.0 & 2{,}353.0 & 2{,}353.0 & 2{,}353.0 & 28.5 \\
HNSC & 10.0 & n.m. & 2{,}231.0 & 2{,}231.0 & 2{,}231.0 & 26.6 \\
KIRC & 10.0 & 2.0 & 2{,}370.0 & 2{,}370.0 & 2{,}370.0 & 11.6 \\
LGG & 10.0 & 2.6 & 2{,}294.0 & 2{,}294.0 & 2{,}294.0 & 34.8 \\
LUAD & 10.0 & 1.0 & 2{,}448.0 & 2{,}448.0 & 2{,}448.0 & 8.2 \\
LUSC & 10.0 & 1.0 & 2{,}442.0 & 2{,}442.0 & 2{,}442.0 & 5.2 \\
OV & 10.0 & n.m. & 2{,}548.0 & 2{,}548.0 & 2{,}548.0 & 15.0 \\
PAAD & 10.0 & 2.0 & 2{,}296.0 & 2{,}296.0 & 2{,}296.0 & 17.0 \\
SKCM & 10.0 & 1.2 & 2{,}325.0 & 2{,}325.0 & 2{,}325.0 & 9.7 \\
STAD & 10.0 & 4.2 & 2{,}681.0 & 2{,}681.0 & 2{,}681.0 & 1.0 \\
\midrule
Mean & 10.0 & 1.8 & 2{,}398.1 & 2{,}398.1 & 2{,}398.1 & 15.9 \\
Cohorts & 11 & 9 & 11 & 11 & 11 & 11 \\
\midrule
METABRIC & 10.0 & n.r. & 1{,}984.0 & 1{,}984.0 & 1{,}984.0 & 176.8 \\
\bottomrule
\end{tabular}

\end{table}

\begin{table}[htbp]
\centering\footnotesize
\caption{Failed folds and, in brackets, null-fit folds per cohort, out of five. A failed
fold produced no prediction. A null-fit fold produced the same predicted survival curve
for every test patient (Section~\ref{S-sec:ml}).}
\label{S-tab:failures}
\setlength{\tabcolsep}{4pt}
\begin{tabular}{@{}l rr rrrr@{}}
\toprule
 & \multicolumn{2}{c}{\pkg{survfcbf}} & \multicolumn{4}{c}{Baselines} \\
\cmidrule(lr){2-3}\cmidrule(l){4-7}
Cohort & top-$K$ & permutation & Variance & CAR & MIM & Lasso Cox \\
\midrule
BLCA & 0 (0) & 5 (0) & 0 (0) & 0 (0) & 0 (0) & 0 (0) \\
BRCA & 0 (0) & 5 (0) & 0 (0) & 0 (0) & 0 (0) & 1 (0) \\
HNSC & 0 (0) & 5 (0) & 0 (0) & 0 (0) & 0 (0) & 0 (0) \\
KIRC & 0 (0) & 4 (0) & 0 (0) & 0 (0) & 0 (0) & 0 (0) \\
LGG & 0 (0) & 1 (0) & 0 (0) & 0 (0) & 0 (0) & 0 (0) \\
LUAD & 0 (0) & 5 (0) & 0 (0) & 0 (1) & 0 (0) & 0 (0) \\
LUSC & 0 (0) & 5 (0) & 0 (0) & 0 (1) & 0 (0) & 1 (0) \\
OV & 0 (0) & 5 (0) & 0 (3) & 0 (0) & 0 (0) & 3 (0) \\
PAAD & 0 (0) & 3 (0) & 0 (0) & 0 (0) & 0 (0) & 0 (0) \\
SKCM & 0 (0) & 4 (0) & 0 (0) & 0 (0) & 0 (0) & 2 (0) \\
STAD & 0 (0) & 0 (0) & 0 (0) & 0 (0) & 0 (0) & 5 (0) \\
\midrule
Total & 0 (0) & 42 (0) & 0 (3) & 0 (2) & 0 (0) & 12 (0) \\
\midrule
METABRIC & 0 (0) & n.r. & 0 (0) & 0 (0) & 0 (0) & 0 (0) \\
\bottomrule
\end{tabular}

\end{table}

\section{Cross-cohort comparison and power}\label{S-sec:bayes}

\paragraph{Bayesian signed-rank test.} The cross-cohort comparisons use the Bayesian
signed-rank test of \citet{benavoli2017}, as implemented in
\code{bayes\_signed\_rank()}: a Dirichlet-process posterior over the per-cohort paired
differences, with a pseudo-observation at zero of prior weight 0.5 and 50{,}000
posterior draws. Each region (\pkg{survfcbf} better, practically equivalent, rival
better) is summarised by the share of draws in which it holds the largest posterior
mass, the summary of the reference implementation, and a region is declared when that
share exceeds 0.95. The region of practical equivalence (ROPE) is $\pm 0.02$ in Uno's C
and $\pm 0.05$ in stability. A cohort enters a comparison only when both selectors have a
value in it (Table~\ref{S-tab:bayes}). The 0.95 rule is applied to \SuppNComparisons{}
comparisons in the tables of this supplement (Tables~\ref{S-tab:rcv}, \ref{S-tab:bayes},
\ref{S-tab:cardinf} and~\ref{S-tab:rcvmatched}), and it does not control an error rate
across them.

\paragraph{Repeated cross-validation.} A single five-fold split leaves resampling noise in
every paired difference. \code{repeated\_cv\_main.R} repeats the split \SuppRcvSplits{}
times per cohort (seeds 1 to 10, seed 1 the published split) for \code{top}-$K$ and the
three rankers on the full feature set, and \code{repeated\_cv.R} does the same at the
matched sizes $K = 5$, 10 and 20 on the protein-coding features, with mRMR and with MIM
and mRMR on the two-bin target (Section~\ref{S-sec:cardinality}). A cohort's score is its
mean over the ten splits, and the Bayesian test runs on those means
(Tables~\ref{S-tab:rcv} and~\ref{S-tab:rcvcohort}). In the headline comparison the gap in
truncated Uno's C is \SuppGapVar{} (95\% $t$ interval \SuppGapLoVar{} to \SuppGapHiVar)
against the variance filter, \SuppGapCar{} (\SuppGapLoCar{} to \SuppGapHiCar) against CAR
and \SuppGapMim{} (\SuppGapLoMim{} to \SuppGapHiMim) against MIM, decided for the ranker in
all three. With the untruncated C the three comparisons stay undecided. METABRIC
(\code{metabric\_repeated\_cv.R}) shows a gap of \SuppMetabricGapMin{} to
\SuppMetabricGapMax{} in every one of its ten splits. The integrated Brier score of
\code{top}-$K$ is \SuppIbsGapMin{} to \SuppIbsGapMax{} higher than that of the rankers, worse
in \SuppIbsWorseMin{} to \SuppIbsWorseMax{} of eleven cohorts but with 95\% $t$ intervals
that include zero, and in METABRIC it is \SuppMetabricIbsMin{} to \SuppMetabricIbsMax{}
higher in all ten splits. On METABRIC's published split a
clinical model of age, Nottingham prognostic index, ER status, HER2 copy-number gain and
PAM50 subtype, fitted with the same harness on the \SuppClinN{} patients with complete
clinical data (\SuppClinEvents{} events, \code{metabric\_clinical.R}), reaches a truncated C
of \SuppClinC, the ten-gene panel alone \SuppClinPanelC{} and both together \SuppClinBothC,
so the panel adds \SuppClinAdd{} (standard error \SuppClinAddSe) to the clinical model.

\begin{table}[htbp]
\centering\footnotesize
\caption{Headline comparison over ten five-fold splits: \pkg{survfcbf} in top-$K$ mode
against the rankers at the ten-percent cut-off, on the cohort means over the splits.
Difference: mean paired difference, \pkg{survfcbf} minus the rival, with its standard error
over the eleven cohorts. Cohorts: cohorts favouring \pkg{survfcbf} and the rival. The next
three columns are the posterior probabilities of the Bayesian signed-rank test, in bold
above 0.95. Splits: splits in which the same comparison, analysed on that split alone,
reaches the same verdict. The integrated Brier score (lower is better) has no region of
practical equivalence and is shown without the test.}
\label{S-tab:rcv}
\setlength{\tabcolsep}{4pt}
\begin{tabular}{@{}l l rrr r rrr r@{}}
\toprule
Measure & Rival & \pkg{survfcbf} & Rival & Difference (SE) & Cohorts & \pkg{survfcbf} & ROPE & Rival & Splits \\
\midrule
Uno's C, truncated & Variance & 0.569 & 0.616 & $-$0.047 (0.012) & 2/9 & 0.00 & 0.02 & \textbf{0.98} & 6/10 \\
 & CAR & 0.569 & 0.613 & $-$0.044 (0.012) & 1/10 & 0.00 & 0.03 & \textbf{0.97} & 7/10 \\
 & MIM & 0.569 & 0.617 & $-$0.048 (0.011) & 0/11 & 0.00 & 0.01 & \textbf{0.99} & 7/10 \\
\addlinespace
Uno's C, untruncated & Variance & 0.561 & 0.598 & $-$0.037 (0.011) & 0/11 & 0.00 & 0.37 & 0.63 & 7/10 \\
 & CAR & 0.561 & 0.607 & $-$0.045 (0.012) & 0/11 & 0.00 & 0.11 & 0.89 & 6/10 \\
 & MIM & 0.561 & 0.604 & $-$0.043 (0.011) & 1/10 & 0.00 & 0.07 & 0.93 & 5/10 \\
\addlinespace
Stability & Variance & 0.109 & 0.942 & $-$0.833 (0.009) & 0/11 & 0.00 & 0.00 & \textbf{1.00} & 10/10 \\
 & CAR & 0.109 & 0.410 & $-$0.301 (0.009) & 0/11 & 0.00 & 0.00 & \textbf{1.00} & 10/10 \\
 & MIM & 0.109 & 0.200 & $-$0.091 (0.018) & 1/10 & 0.00 & 0.04 & \textbf{0.96} & 7/10 \\
\addlinespace
Integrated Brier score & Variance & 0.230 & 0.229 & 0.001 (0.008) & 3/8 &  &  &  &  \\
 & CAR & 0.230 & 0.228 & 0.002 (0.008) & 2/9 &  &  &  &  \\
 & MIM & 0.230 & 0.226 & 0.004 (0.007) & 1/10 &  &  &  &  \\
\bottomrule
\end{tabular}

\end{table}

\begin{table}[htbp]
\centering\footnotesize
\caption{Truncated Uno's C per cohort over ten five-fold splits: mean (standard deviation
over the splits). METABRIC ran on a workstation (Section~\ref{S-sec:software}).}
\label{S-tab:rcvcohort}
\begin{tabular}{@{}l rrrr@{}}
\toprule
Cohort & top-$K$ & Variance & CAR & MIM \\
\midrule
BLCA & 0.536 (0.021) & 0.640 (0.021) & 0.654 (0.017) & 0.654 (0.012) \\
BRCA & 0.503 (0.034) & 0.620 (0.025) & 0.602 (0.031) & 0.604 (0.028) \\
HNSC & 0.587 (0.024) & 0.633 (0.019) & 0.644 (0.021) & 0.650 (0.017) \\
KIRC & 0.586 (0.031) & 0.653 (0.015) & 0.650 (0.012) & 0.633 (0.017) \\
LGG & 0.746 (0.027) & 0.816 (0.017) & 0.792 (0.014) & 0.795 (0.017) \\
LUAD & 0.564 (0.029) & 0.590 (0.016) & 0.579 (0.014) & 0.585 (0.024) \\
LUSC & 0.500 (0.039) & 0.531 (0.017) & 0.511 (0.017) & 0.509 (0.022) \\
OV & 0.502 (0.025) & 0.527 (0.026) & 0.534 (0.018) & 0.530 (0.026) \\
PAAD & 0.605 (0.031) & 0.657 (0.017) & 0.637 (0.024) & 0.659 (0.019) \\
SKCM & 0.579 (0.030) & 0.570 (0.015) & 0.570 (0.018) & 0.583 (0.017) \\
STAD & 0.548 (0.037) & 0.541 (0.035) & 0.571 (0.035) & 0.582 (0.034) \\
\midrule
Mean & 0.569 & 0.616 & 0.613 & 0.617 \\
\midrule
METABRIC & 0.583 (0.005) & 0.632 (0.002) & 0.627 (0.005) & 0.630 (0.003) \\
\bottomrule
\end{tabular}

\end{table}

\paragraph{Summary of the posterior.} The posterior can also be summarised by the mean
mass of each region over the draws, which is the default of some implementations. Under
that summary the probability that the rival is better in the headline comparison is
\SuppMeanVar{} for the variance filter, \SuppMeanCar{} for CAR and \SuppMeanMim{} for MIM,
so \SuppMeanDecidedText{} at 0.95. The $t$ intervals above do not depend on that choice,
and none of them includes zero.

\paragraph{Width of the ROPE.} Table~\ref{S-tab:rope} repeats the comparisons with ROPEs of
$\pm 0.01$ and $\pm 0.03$ in Uno's C and $\pm 0.03$ and $\pm 0.07$ in stability. The
headline comparison is decided at $\pm 0.01$ and $\pm 0.02$ and not at $\pm 0.03$. At every
matched size the practical equivalence with mRMR holds at $\pm 0.02$ and $\pm 0.03$ (the
table shows ten genes).

\begin{table}[htbp]
\centering\footnotesize
\caption{Verdict of the 0.95 rule by width of the ROPE, on the cohort means over ten
splits: \pkg{survfcbf} better, rival better, equivalent (practical equivalence) or
undecided. ``Mean'': the mean-mass summary of the posterior instead of the argmax.}
\label{S-tab:rope}
\setlength{\tabcolsep}{4pt}
\begin{tabular}{@{}l llll lll@{}}
\toprule
 & \multicolumn{4}{c}{Uno's C, truncated} & \multicolumn{3}{c}{Stability} \\
\cmidrule(lr){2-5}\cmidrule(l){6-8}
Rival & $\pm 0.01$ & $\pm 0.02$ & $\pm 0.03$ & $\pm 0.02$, mean & $\pm 0.03$ & $\pm 0.05$ & $\pm 0.07$ \\
\midrule
\multicolumn{8}{@{}l}{\textit{A. Headline, at the rankers' own panel size}} \\
Variance & rival & rival & undecided & undecided & rival & rival & rival \\
CAR & rival & rival & undecided & undecided & rival & rival & rival \\
MIM & rival & rival & undecided & undecided & rival & rival & undecided \\
\midrule
\multicolumn{8}{@{}l}{\textit{B. Matched panel size, $K = 10$}} \\
Variance & undecided & undecided & undecided & undecided & rival & rival & rival \\
CAR & undecided & undecided & equivalent & undecided & rival & rival & rival \\
MIM & undecided & undecided & equivalent & undecided & \pkg{survfcbf} & undecided & undecided \\
mRMR & undecided & equivalent & equivalent & undecided & \pkg{survfcbf} & \pkg{survfcbf} & undecided \\
MIM, $q = 2$ & undecided & equivalent & equivalent & undecided & undecided & undecided & equivalent \\
mRMR, $q = 2$ & equivalent & equivalent & equivalent & equivalent & equivalent & equivalent & equivalent \\
\bottomrule
\end{tabular}

\end{table}

\begin{table}[htbp]
\centering\footnotesize
\caption{Bayesian signed-rank comparison of \pkg{survfcbf} in top-$K$ mode with every other
selector across the TCGA cohorts, on the published split with the untruncated Uno's C.
Each probability is the share of posterior draws in which the region holds the largest
mass, as written by \code{analyze\_cross\_dataset.R}, rounded to two decimals.
Probabilities above 0.95, the decision threshold, are in bold.}
\label{S-tab:bayes}
\begin{tabular}{@{}l rrrr rrrr@{}}
\toprule
 & \multicolumn{4}{c}{Uno's C, ROPE $\pm 0.02$} & \multicolumn{4}{c}{Stability, ROPE $\pm 0.05$} \\
\cmidrule(lr){2-5}\cmidrule(l){6-9}
Rival & Cohorts & top-$K$ & ROPE & Rival & Cohorts & top-$K$ & ROPE & Rival \\
\midrule
Variance & 11 & 0.00 & 0.00 & \textbf{1.00} & 11 & 0.00 & 0.00 & \textbf{1.00} \\
CAR & 11 & 0.00 & 0.00 & \textbf{1.00} & 11 & 0.00 & 0.00 & \textbf{1.00} \\
MIM & 11 & 0.00 & 0.00 & \textbf{1.00} & 11 & 0.00 & 0.02 & \textbf{0.98} \\
Lasso Cox & 10 & 0.00 & 0.01 & \textbf{0.99} & 11 & 0.00 & 0.07 & 0.93 \\
\pkg{survfcbf}, permutation & 5 & 0.02 & 0.53 & 0.45 & 8 & \textbf{0.98} & 0.02 & 0.00 \\
\bottomrule
\end{tabular}

\end{table}

\paragraph{Power.} How many cohorts a comparison needs depends on how much the paired
difference between two selectors varies from cohort to cohort (\code{power\_analysis.R}).
That standard deviation was estimated on the cohort means of the truncated C over ten
splits, separately for the headline comparison (\code{top}-$K$ against the three rankers,
\SuppPairsHead{} pairs, from \SuppSigmaHeadMin{} to \SuppSigmaHeadMax, median
\SuppSigmaHead) and for the matched sizes (the same rankers at $K = 5$, 10 and 20,
\SuppPairsMatched{} pairs, median \SuppSigmaMatched, from \SuppSigmaMatchedMin{} to
\SuppSigmaMatchedMax, the largest \SuppSigmaMatchedMaxWhat). A median over every pair of
selectors would depend on which variants a run contains, since near-identical variants
give very small standard deviations, so the other pairs enter no estimate. Panel A of
Table~\ref{S-tab:power} gives the cohorts a two-sided paired test at level 0.05 needs to
detect a difference, from the normal approximation
$N = (z_{1-\alpha/2} + z_{1-\beta})^2 \sigma^2 / \Delta^2$
\citep{boulesteix2017}. Eleven cohorts give such a test 80\% power for a difference of
\SuppDeltaHead{} in the headline comparison, and detecting 0.03 needs
\SuppNeedThreeHead{} cohorts there, \SuppNeedThreeMatched{} at the matched sizes and
\SuppNeedThreeMatchedMax{} at their largest standard deviation. Panel B gives the power of
the 0.95 rule of the Bayesian test at each median, from \SuppPowerReps{} simulated
benchmarks per cell with normally distributed paired differences, each test run with
1{,}500 posterior draws. With eleven cohorts it declares the better selector in a share of
\SuppPowerHeadThree{} of the simulated benchmarks at a headline difference of 0.03 and of
\SuppPowerHeadFive{} at 0.05. On the published split alone, with the untruncated C, the
standard deviations of the paired differences have a median of \SuppSdLegacy.

\begin{table}[htbp]
\centering\footnotesize
\caption{Number of cohorts needed to resolve a difference in truncated Uno's C, for the
headline comparison and for the matched sizes. Panel A: cohorts a paired test needs at
80\% and 90\% power at the median standard deviation of paired differences, and at 80\%
power at the largest one. Panel B: share of simulated benchmarks in which the Bayesian
signed-rank test declares the better selector, with its Monte Carlo standard error in
brackets.}
\label{S-tab:power}
\begin{tabular}{@{}l rrr rrr@{}}
\toprule
\multicolumn{7}{@{}l}{\textit{A. Cohorts needed by a paired test (two-sided $\alpha = 0.05$)}} \\
 & \multicolumn{3}{c}{Headline, SD 0.038} & \multicolumn{3}{c}{Matched size, SD 0.030} \\
\cmidrule(lr){2-4}\cmidrule(l){5-7}
Difference & 80\% & 90\% & 80\%, largest SD & 80\% & 90\% & 80\%, largest SD \\
\midrule
0.01 & 115 & 154 & 129 & 70 & 94 & 475 \\
0.02 & 29 & 39 & 33 & 18 & 24 & 119 \\
0.03 & 13 & 18 & 15 & 8 & 11 & 53 \\
0.05 & 5 & 7 & 6 & 3 & 4 & 19 \\
0.10 & 2 & 2 & 2 & 1 & 1 & 5 \\
\midrule
\multicolumn{7}{@{}l}{\textit{B. Power of the 0.95 rule (Bayesian signed-rank, ROPE $\pm 0.02$)}} \\
Difference & 11 cohorts & 30 cohorts & 100 cohorts & 11 cohorts & 30 cohorts & 100 cohorts \\
\midrule
0.03 & 0.21 (0.02) & 0.49 (0.02) & 0.92 (0.01) & 0.28 (0.02) & 0.57 (0.02) & 0.97 (0.01) \\
0.05 & 0.79 (0.02) & 1.00 (0.00) & 1.00 (0.00) & 0.91 (0.01) & 1.00 (0.00) & 1.00 (0.00) \\
0.08 & 1.00 (0.00) & 1.00 (0.00) & 1.00 (0.00) & 1.00 (0.00) & 1.00 (0.00) & 1.00 (0.00) \\
0.12 & 1.00 (0.00) & 1.00 (0.00) & 1.00 (0.00) & 1.00 (0.00) & 1.00 (0.00) & 1.00 (0.00) \\
\bottomrule
\end{tabular}

\end{table}

\section{Matched panel size}\label{S-sec:cardinality}

Comparing selectors at their own panel sizes mixes the effect of the selection rule with
that of the panel size. To separate them, \pkg{survfcbf} and the three rankers were cut
to a panel size $K \in \{5, 10, 20, 50, 100, 500, 2000\}$ and run through the same
harness, with the same
five outer folds per cohort, the same seed and the same downstream model
(\code{cardinality\_curve.R}). The rankers keep their first $K$ features. \pkg{survfcbf}
runs in top-$K$ mode at $q = 2$ up to $K = 100$, a bound set to limit its cost. This
analysis uses the protein-coding RNA features only (\SuppCodingMin{} to \SuppCodingMax{}
per cohort, Table~\ref{S-tab:cohorts}). No fold failed at any size.

\paragraph{The size of a top-$K$ panel.} In top-$K$ mode the relevance threshold is zero,
so the Markov-blanket pruning runs over every feature, and it stops early once $K$
predominant features remain. When the whole feature set holds fewer than $K$ predominant
features, the panel is smaller than $K$. At $q = 2$ this happens in
\SuppCardShortTwenty{} cohorts at $K = 20$ and in every cohort from $K = 50$, where the
complete set of predominant features holds \SuppCardPanelMin{} to \SuppCardPanelMax{}
genes (mean \SuppCardPanelMean). The panel size is therefore matched exactly up to
$K = \SuppCardExactMax$, nearly at $K = 20$, and not beyond: the rows with $K \ge 50$
compare panels of about \SuppCardPanelMean{} genes with ranker panels of $K$ genes
(Table~\ref{S-tab:cardinality}, panel A).

Table~\ref{S-tab:cardinality} gives the means over cohorts on the published split, with
the untruncated C, and the number of null fits, folds in which the inner cross-validation
of the downstream model found no penalty that improved on the null model.
Table~\ref{S-tab:cardinf} gives the paired comparisons of that split, with the Bayesian
signed-rank test of Section~\ref{S-sec:bayes} (ROPE $\pm 0.02$ in Uno's C). The comparisons
at the matched sizes $K = 5$, 10 and 20 were repeated over ten splits with the truncated C
(\code{repeated\_cv.R}, Table~\ref{S-tab:rcvmatched}). There \pkg{survfcbf} is
\SuppRcvVarLo{} to \SuppRcvVarHi{} ahead of the variance filter and \SuppRcvCarLo{} to
\SuppRcvCarHi{} ahead of CAR, and \SuppRcvMimLo{} to \SuppRcvMimHi{} behind MIM, with no
comparison decided except practical equivalence with CAR at $K = 20$. The null fits of the
variance filter (\SuppRcvNullVarLo\% to \SuppRcvNullVarHi\% of folds at these sizes,
against \SuppRcvNullTopk{} fold in \SuppRcvNullTopkFolds{} for \pkg{survfcbf}) pull its mean
towards 0.5 and are part of what puts it behind.

\begin{table}[htbp]
\centering\footnotesize
\caption{Selectors cut to a panel size $K$, over the eleven TCGA cohorts and their
protein-coding features. The rankers keep exactly $K$ features. \pkg{survfcbf} keeps
fewer when the feature set holds fewer than $K$ predominant features (panel A, mean and
range over cohorts of the mean panel per fold). n.r., not run.}
\label{S-tab:cardinality}
\begin{tabular}{@{}r lrrr@{}}
\toprule
$K$ & \pkg{survfcbf} & Variance & CAR & MIM \\
\midrule
\multicolumn{5}{@{}l}{\textit{A. Panel size, mean (range) over cohorts}} \\
5 & 5 & 5 & 5 & 5 \\
10 & 10 & 10 & 10 & 10 \\
20 & 19.5 (17.2--20.0) & 20 & 20 & 20 \\
50 & 26.6 (17.6--43.2) & 50 & 50 & 50 \\
100 & 26.7 (17.6--43.8) & 100 & 100 & 100 \\
500 & n.r. & 500 & 500 & 500 \\
2{,}000 & n.r. & 2{,}000 & 2{,}000 & 2{,}000 \\
\midrule
\multicolumn{5}{@{}l}{\textit{B. Uno's C, mean over cohorts}} \\
5 & 0.542 & 0.524 & 0.530 & 0.578 \\
10 & 0.548 & 0.523 & 0.544 & 0.582 \\
20 & 0.566 & 0.535 & 0.561 & 0.588 \\
50 & 0.576 & 0.547 & 0.583 & 0.604 \\
100 & 0.575 & 0.564 & 0.596 & 0.600 \\
500 & n.r. & 0.598 & 0.603 & 0.606 \\
2{,}000 & n.r. & 0.603 & 0.606 & 0.602 \\
\midrule
\multicolumn{5}{@{}l}{\textit{C. Stability, mean over cohorts}} \\
5 & 0.076 & 0.884 & 0.243 & 0.040 \\
10 & 0.090 & 0.905 & 0.269 & 0.057 \\
20 & 0.102 & 0.921 & 0.295 & 0.063 \\
50 & 0.101 & 0.920 & 0.313 & 0.087 \\
100 & 0.101 & 0.923 & 0.328 & 0.104 \\
500 & n.r. & 0.932 & 0.372 & 0.155 \\
2{,}000 & n.r. & 0.951 & 0.404 & 0.211 \\
\midrule
\multicolumn{5}{@{}l}{\textit{D. Null-fit folds, out of 55}} \\
5 & 0 & 13 & 7 & 1 \\
10 & 0 & 10 & 5 & 1 \\
20 & 0 & 8 & 2 & 1 \\
50 & 0 & 2 & 3 & 0 \\
100 & 0 & 5 & 0 & 0 \\
500 & n.r. & 3 & 0 & 0 \\
2{,}000 & n.r. & 3 & 2 & 0 \\
\bottomrule
\end{tabular}

\end{table}

\begin{table}[htbp]
\centering\footnotesize
\caption{Paired comparison of \pkg{survfcbf} with each ranker cut to a panel size $K$, on
the published split with the untruncated Uno's C.
\pkg{survfcbf} size: its mean panel size, which falls short of $K$ from $K = 20$
(Section~\ref{S-sec:cardinality}). Difference: mean paired difference in Uno's C,
\pkg{survfcbf} minus the rival, with its standard error over cohorts. Cohorts: cohorts favouring \pkg{survfcbf} and the rival. The
last three columns are the posterior probabilities of the Bayesian signed-rank test. No
region exceeds 0.95, so no comparison is decided with eleven cohorts.}
\label{S-tab:cardinf}
\begin{tabular}{@{}l r r r r rrr@{}}
\toprule
Rival & $K$ & \pkg{survfcbf} size & Difference (SE) & Cohorts & \pkg{survfcbf} & ROPE & Rival \\
\midrule
Variance & 5 & 5.0 & 0.018 (0.025) & 6/5 & 0.52 & 0.34 & 0.14 \\
 & 10 & 10.0 & 0.025 (0.019) & 7/4 & 0.81 & 0.12 & 0.07 \\
 & 20 & 19.5 & 0.030 (0.010) & 9/2 & 0.90 & 0.10 & 0.00 \\
 & 50 & 26.6 & 0.029 (0.015) & 9/2 & 0.87 & 0.12 & 0.01 \\
 & 100 & 26.7 & 0.011 (0.018) & 5/6 & 0.55 & 0.32 & 0.13 \\
\addlinespace
CAR & 5 & 5.0 & 0.012 (0.015) & 6/5 & 0.46 & 0.50 & 0.04 \\
 & 10 & 10.0 & 0.004 (0.012) & 5/6 & 0.15 & 0.81 & 0.04 \\
 & 20 & 19.5 & 0.005 (0.010) & 6/5 & 0.12 & 0.85 & 0.03 \\
 & 50 & 26.6 & $-$0.007 (0.026) & 4/7 & 0.32 & 0.07 & 0.61 \\
 & 100 & 26.7 & $-$0.020 (0.018) & 2/9 & 0.11 & 0.23 & 0.65 \\
\addlinespace
MIM & 5 & 5.0 & $-$0.037 (0.019) & 3/8 & 0.01 & 0.35 & 0.65 \\
 & 10 & 10.0 & $-$0.034 (0.023) & 4/7 & 0.07 & 0.18 & 0.76 \\
 & 20 & 19.5 & $-$0.022 (0.020) & 4/7 & 0.06 & 0.27 & 0.67 \\
 & 50 & 26.6 & $-$0.028 (0.016) & 3/8 & 0.00 & 0.22 & 0.78 \\
 & 100 & 26.7 & $-$0.025 (0.014) & 2/9 & 0.00 & 0.47 & 0.53 \\
\bottomrule
\end{tabular}

\end{table}

\begin{table}[htbp]
\centering\footnotesize
\caption{Paired comparison of \pkg{survfcbf} (top-$K$, $q = 2$) with each rival cut to the
same panel size $K$, on the cohort means over ten five-fold splits. Difference: mean paired
difference, \pkg{survfcbf} minus the rival, with its standard error over the eleven cohorts.
Cohorts: cohorts favouring \pkg{survfcbf} and the rival. Verdict: the 0.95 rule of the
Bayesian signed-rank test, with a ROPE of $\pm 0.02$ in Uno's C and $\pm 0.05$ in
stability.}
\label{S-tab:rcvmatched}
\setlength{\tabcolsep}{4pt}
\begin{tabular}{@{}l r r r l r r l@{}}
\toprule
 &  & \multicolumn{3}{c}{Uno's C, truncated} & \multicolumn{3}{c}{Stability} \\
\cmidrule(lr){3-5}\cmidrule(l){6-8}
Rival & $K$ & Difference (SE) & Cohorts & Verdict & Difference (SE) & Cohorts & Verdict \\
\midrule
Variance & 5 & 0.022 (0.023) & 8/3 & undecided & $-$0.779 (0.033) & 0/11 & rival \\
 & 10 & 0.029 (0.018) & 8/3 & undecided & $-$0.803 (0.015) & 0/11 & rival \\
 & 20 & 0.017 (0.015) & 6/5 & undecided & $-$0.803 (0.009) & 0/11 & rival \\
\addlinespace
CAR & 5 & 0.021 (0.014) & 7/4 & undecided & $-$0.166 (0.027) & 0/11 & rival \\
 & 10 & 0.012 (0.008) & 8/3 & undecided & $-$0.174 (0.019) & 0/11 & rival \\
 & 20 & 0.008 (0.007) & 5/6 & equivalent & $-$0.187 (0.016) & 0/11 & rival \\
\addlinespace
MIM & 5 & $-$0.013 (0.009) & 4/7 & undecided & 0.064 (0.012) & 9/2 & undecided \\
 & 10 & $-$0.014 (0.008) & 4/7 & undecided & 0.062 (0.015) & 10/1 & undecided \\
 & 20 & $-$0.014 (0.009) & 3/8 & undecided & 0.053 (0.016) & 9/2 & undecided \\
\addlinespace
mRMR & 5 & $-$0.002 (0.006) & 6/5 & equivalent & 0.082 (0.010) & 11/0 & \pkg{survfcbf} \\
 & 10 & $-$0.003 (0.007) & 6/5 & equivalent & 0.084 (0.009) & 11/0 & \pkg{survfcbf} \\
 & 20 & $-$0.006 (0.007) & 5/6 & equivalent & 0.077 (0.008) & 11/0 & \pkg{survfcbf} \\
\addlinespace
MIM, $q = 2$ & 5 & $-$0.009 (0.004) & 3/8 & equivalent & $-$0.033 (0.013) & 3/8 & undecided \\
 & 10 & $-$0.006 (0.004) & 4/7 & equivalent & $-$0.035 (0.010) & 2/9 & undecided \\
 & 20 & $-$0.003 (0.005) & 5/6 & equivalent & $-$0.045 (0.009) & 0/11 & undecided \\
\addlinespace
mRMR, $q = 2$ & 5 & $-$0.002 (0.004) & 5/6 & equivalent & 0.014 (0.008) & 8/3 & equivalent \\
 & 10 & $-$0.002 (0.003) & 5/6 & equivalent & $-$0.003 (0.005) & 6/5 & equivalent \\
 & 20 & $-$0.006 (0.004) & 3/8 & equivalent & $-$0.038 (0.006) & 0/11 & equivalent \\
\bottomrule
\end{tabular}

\end{table}

\paragraph{Redundancy within the panel and mRMR.} \code{panel\_redundancy.R} reruns the
sizes $K = 5$, 10 and 20 on the same folds and adds mRMR (\code{MRMR} of \pkg{praznik} on
the discretised martingale residuals that MIM uses, \citealp{kursa2021}, keeping the first
$K$ features of its greedy order), MIM and mRMR on the two-bin target of \code{top}-$K$
($q = 2$), and top-$K$ at $q = 4$ and $q = 10$ for $K = 10$. For each fold it records two
measures of redundancy on the training rows: the mean absolute Pearson correlation
between the selected genes, and the share of gene pairs with $|r| \ge 0.5$. Its
untruncated C and stability for \pkg{survfcbf} and the three rankers reproduce
Table~\ref{S-tab:cardinality} exactly. Table~\ref{S-tab:redundancy} gives the means over
cohorts. At $K = 10$ Markov-blanket pruning halves the redundancy of MIM on the same
two-bin target, which scores the same relevance without pruning (mean $|r|$ \SuppRedTopk{}
against \SuppRedMimQtwo, lower in \SuppRedBelowMimQtwo{} of eleven cohorts, and
\SuppRedMim{} for MIM at its default binning), and the share of strongly correlated pairs
falls from \SuppShareMim\% to \SuppShareTopk\%. mRMR (\SuppRedMrmr, and \SuppRedMrmrQtwo{}
at $q = 2$) and CAR (\SuppRedCar) reach a similar or lower level, the first by penalising
redundancy and the second by adjusting for the correlation between features, and mRMR at
$q = 2$ is below \pkg{survfcbf} in all but \SuppRedBelowMrmrQtwo{} cohorts. Over ten splits
(Table~\ref{S-tab:rcvmatched}), \pkg{survfcbf} and mRMR are practically equivalent in
truncated Uno's C at every size, at the default binning and at $q = 2$. \pkg{survfcbf} is
the more stable against mRMR at its default binning, decided at every size, and
practically equivalent to it at $q = 2$.

\paragraph{What the trace records.} With the panel, \pkg{survfcbf} returns every gene it
pruned and the panel gene that pruned it. \code{panel\_groups.R} reads these groups for the
ten-gene panels of the eleven cohorts, which it reproduces in all \SuppGrpFolds{} folds. At
$q = 2$ the relevance of a gene is small, so the pruning condition holds at modest
correlation. Across cohorts, the median absolute correlation of a pruned gene with its
panel gene ranged from \SuppGrpRMin{} to \SuppGrpRMax, against a mean of \SuppGrpBgMin{}
to \SuppGrpBgMax{} between a panel gene and all other genes, and \SuppGrpHalfMin\% to
\SuppGrpHalfMax\% of pruned genes reached 0.5. In every cohort the median pruned gene was
among the \SuppGrpTopPct\% of genes most correlated with its panel gene. FCBF assigns each
pruned gene to the most relevant panel member that satisfies the condition, not to the
closest one, and that member was also the closest in \SuppGrpClosestMin\% to
\SuppGrpClosestMax\% of cases. Of the panel genes, \SuppGrpEmpty\% pruned nothing, and the
largest group held \SuppGrpMaxGroup{} genes. The trace therefore marks neighbourhoods of
the panel genes rather than near-copies.

\paragraph{The binning of the target.} Against the default $q = 2$, top-$K$ at $q = 4$ and
$q = 10$ changes the mean truncated Uno's C of the published split by \SuppQfourC{} and
\SuppQtenC, with per-cohort differences from \SuppQdiffMin{} to \SuppQdiffMax, and lowers
stability in \SuppQfourLessStable{} and \SuppQtenLessStable{} of eleven cohorts. At
$q = 10$ the panel is also more redundant in \SuppQtenMoreRedundant{} cohorts, and both
finer binnings give panels that predict censoring beyond the 95th percentile of random
panels in \SuppCensAboveQfour{} cohorts (Section~\ref{S-sec:selectors}). The default
$q = 2$ is kept for its stability, its lower redundancy and because its panels carry no
follow-up signal beyond that of random genes, not for discrimination.

\paragraph{The similarity threshold of the stability index.} The index counts two genes as
similar at $|r| \ge 0.9$. At 0.7 and 0.5 (\code{stability\_theta.R}, the threshold 0.5
with \SuppThetaDrawsLow{} draws for the correction for chance instead of
\SuppThetaDraws) the ordering does not change (Table~\ref{S-tab:theta}). \pkg{survfcbf}
stays below MIM on the two-bin target in \SuppThetaLowerMin{} or \SuppThetaLowerMax{} of
eleven cohorts, and the gap widens as the threshold falls (\SuppThetaGapNine,
\SuppThetaGapSeven, \SuppThetaGapFive), because a pruned panel holds no near-duplicates
for a looser threshold to match. At 0.9 the run reproduces Table~\ref{S-tab:redundancy}
exactly.

\begin{table}[htbp]
\centering\footnotesize
\caption{Matched panel sizes with mRMR and with MIM and mRMR on the two-bin target, and
the redundancy within each panel, as means over the eleven TCGA cohorts and their
protein-coding features, on the published split. \pkg{survfcbf}: top-$K$ at $q = 2$ unless
stated. Uno's C: truncated. Mean $|r|$: mean absolute Pearson correlation between the
selected genes. $|r| \ge 0.5$: share of gene pairs at or above 0.5. Both are computed on
the training rows of each fold. Null fits: folds, out of 55, in which the downstream model
fell back to the null model.}
\label{S-tab:redundancy}
\begin{tabular}{@{}l r rrrrr@{}}
\toprule
Selector & $K$ & Uno's C & Stability & Mean $|r|$ & $|r| \ge 0.5$ (\%) & Null fits \\
\midrule
\pkg{survfcbf} & 5 & 0.555 & 0.076 & 0.161 & 1.6 & 0 \\
mRMR & 5 & 0.553 & 0.005 & 0.109 & 0.2 & 4 \\
mRMR, $q = 2$ & 5 & 0.549 & 0.065 & 0.111 & 0.4 & 0 \\
MIM, $q = 2$ & 5 & 0.561 & 0.125 & 0.271 & 16.2 & 2 \\
Variance & 5 & 0.526 & 0.884 & 0.501 & 45.1 & 13 \\
CAR & 5 & 0.522 & 0.243 & 0.107 & 1.8 & 7 \\
MIM & 5 & 0.572 & 0.040 & 0.272 & 15.5 & 1 \\
\addlinespace
\pkg{survfcbf} & 10 & 0.572 & 0.090 & 0.134 & 0.8 & 0 \\
mRMR & 10 & 0.570 & 0.022 & 0.119 & 0.6 & 3 \\
mRMR, $q = 2$ & 10 & 0.565 & 0.108 & 0.121 & 0.8 & 0 \\
MIM, $q = 2$ & 10 & 0.578 & 0.138 & 0.261 & 12.7 & 1 \\
Variance & 10 & 0.543 & 0.905 & 0.432 & 38.8 & 10 \\
CAR & 10 & 0.550 & 0.269 & 0.105 & 1.2 & 5 \\
MIM & 10 & 0.582 & 0.057 & 0.266 & 15.3 & 1 \\
\pkg{survfcbf}, $q = 4$ & 10 & 0.585 & 0.041 & 0.149 & 0.9 & 0 \\
\pkg{survfcbf}, $q = 10$ & 10 & 0.560 & 0.044 & 0.165 & 1.7 & 2 \\
\addlinespace
\pkg{survfcbf} & 20 & 0.583 & 0.102 & 0.109 & 0.4 & 0 \\
mRMR & 20 & 0.570 & 0.042 & 0.127 & 0.8 & 0 \\
mRMR, $q = 2$ & 20 & 0.570 & 0.147 & 0.125 & 0.7 & 0 \\
MIM, $q = 2$ & 20 & 0.582 & 0.151 & 0.253 & 11.6 & 0 \\
Variance & 20 & 0.566 & 0.921 & 0.329 & 26.9 & 8 \\
CAR & 20 & 0.580 & 0.295 & 0.105 & 1.0 & 2 \\
MIM & 20 & 0.584 & 0.063 & 0.262 & 14.7 & 1 \\
\bottomrule
\end{tabular}

\end{table}

\begin{table}[htbp]
\centering\footnotesize
\caption{Selection stability at ten genes by the similarity threshold $\theta$ of the
index, as means over the eleven TCGA cohorts on the published split. At $\theta = 0.5$ the
correction for chance uses \SuppThetaDrawsLow{} draws instead of \SuppThetaDraws, so those
values are approximate.}
\label{S-tab:theta}
\begin{tabular}{@{}l rrr@{}}
\toprule
Selector, $K = 10$ & $\theta = 0.9$ & $\theta = 0.7$ & $\theta = 0.5$ \\
\midrule
\pkg{survfcbf} & 0.090 & 0.093 & 0.050 \\
mRMR & 0.022 & 0.024 & 0.009 \\
mRMR, $q = 2$ & 0.108 & 0.113 & 0.064 \\
MIM, $q = 2$ & 0.138 & 0.160 & 0.178 \\
Variance & 0.905 & 0.907 & 0.902 \\
CAR & 0.269 & 0.263 & 0.184 \\
MIM & 0.057 & 0.127 & 0.290 \\
\midrule
\pkg{survfcbf} minus MIM, $q = 2$ & $-$0.047 & $-$0.066 & $-$0.128 \\
Cohorts with \pkg{survfcbf} lower & 10 & 10 & 9 \\
\bottomrule
\end{tabular}

\end{table}

\section{Gene-set enrichment}\label{S-sec:gsea}

For each TCGA cohort, the per-gene score of every ranker was computed once on the whole
cohort over its protein-coding features, defined by the HGNC list of protein-coding genes
\citep{seal2023} (\SuppCodingIds{} Ensembl identifiers, snapshot of \SuppCodingDate). The
score of \pkg{survfcbf} is the relevance of its top-$K$ mode, the SU of each gene with
the martingale residuals at $q = 2$, with ties ordered by the same Spearman key as the
selector (Section~\ref{S-sec:selectors}). The baselines use their own scores: the
variance, the absolute CAR score and the mutual information. Genes were mapped from
Ensembl identifiers to symbols with \pkg{org.Hs.eg.db}. In each cohort
\SuppNoSymbolMin{} to \SuppNoSymbolMax{} genes have no symbol and \SuppDupMin{} to
\SuppDupMax{} map to a symbol already taken, and both are dropped
(Table~\ref{S-tab:gsea}). The remaining genes were ranked by score and tested with
\pkg{fgsea} \citep{korotkevich2016}
against the Hallmark \citep{liberzon2015} and oncogenic signature (C6)
\citep{liberzon2011} collections of MSigDB, \SuppSetsTotal{} sets in all
(\SuppSetsHallmark{} Hallmark and \SuppSetsCsix{} C6, \pkg{msigdbr} \SuppMsigdbrVersion,
snapshot of \SuppMsigdbDate). The test is one-sided (\code{scoreType = "pos"}), so a set
is enriched when its genes concentrate at the top of the ranking, and only sets of 5 to
500 genes are tested. $P$ values are adjusted by the Benjamini-Hochberg procedure within
each cohort and ranking. The analysis describes the scores of whole cohorts and makes no
prediction, so it is not cross-validated.

A preranked test of this kind asks whether a set's genes sit high in a ranking, and gene
sets are groups of co-expressed genes, so any ranking that follows co-expression enriches
them. The variance score does not use the outcome, and its ranking is enriched in
\SuppGseaVariance{} of eleven cohorts, more than any ranking that does. Enrichment of the
\pkg{survfcbf} ranking, in \SuppGseaSurvfcbf{} of eleven cohorts, therefore shows that its
relevance follows co-expressed programmes, not that those programmes carry survival
signal. Testing that would need a null that permutes the outcome rather than the genes.

\begin{table}[htbp]
\centering\footnotesize
\caption{Gene sets with a Benjamini-Hochberg adjusted $P < 0.05$, per cohort and ranking.
Genes: protein-coding RNA features of the cohort, those lost in the mapping to symbols
(no symbol or a duplicated one) and those ranked. Sets tested: sets of 5 to 500 genes
after mapping to the cohort's genes.}
\label{S-tab:gsea}
\begin{tabular}{@{}l rrr r rrrr@{}}
\toprule
 & \multicolumn{3}{c}{Genes} &  & \multicolumn{4}{c}{Significant sets} \\
\cmidrule(lr){2-4}\cmidrule(l){6-9}
Cohort & Coding & Lost & Ranked & Sets tested & \pkg{survfcbf} & Variance & CAR & MIM \\
\midrule
BLCA & 16{,}390 & 27 & 16{,}363 & 239 & 4 & 73 & 0 & 2 \\
BRCA & 16{,}307 & 24 & 16{,}283 & 239 & 32 & 61 & 11 & 3 \\
HNSC & 16{,}202 & 26 & 16{,}176 & 239 & 0 & 50 & 2 & 1 \\
KIRC & 16{,}206 & 25 & 16{,}181 & 239 & 0 & 52 & 3 & 5 \\
LGG & 15{,}993 & 23 & 15{,}970 & 239 & 2 & 63 & 0 & 15 \\
LUAD & 16{,}527 & 26 & 16{,}501 & 239 & 18 & 50 & 0 & 1 \\
LUSC & 16{,}595 & 27 & 16{,}568 & 239 & 0 & 49 & 0 & 0 \\
OV & 16{,}492 & 28 & 16{,}464 & 239 & 0 & 66 & 0 & 0 \\
PAAD & 16{,}386 & 24 & 16{,}362 & 239 & 4 & 45 & 5 & 2 \\
SKCM & 16{,}115 & 24 & 16{,}091 & 239 & 14 & 75 & 0 & 1 \\
STAD & 16{,}899 & 28 & 16{,}871 & 239 & 12 & 45 & 1 & 13 \\
\midrule
Cohorts with a set &  &  &  &  & 7 & 11 & 5 & 9 \\
\bottomrule
\end{tabular}

\end{table}

\section{Simulation}\label{S-sec:simulation}

\paragraph{Design.} The simulation grid
(tier \code{cesga\_core} of \code{grid\_lib.R}) varies one factor at a
time around a reference configuration with $n = 600$, 50\% censoring and $p = 2{,}000$:
the sample size ($n$ from 150 to 2{,}400), the censored fraction (20\%, 50\% and 70\%)
and the number of features ($p$ = 200, 1{,}000 and 5{,}000). The reference configuration
belongs to both the sample size and the censoring axes and is counted once, which gives
\SuppSimConfigs{} distinct configurations with \SuppSimReps{} replicates each,
\SuppSimDatasets{} datasets in all. Each dataset has three driver features, each the first
member of a cluster of six features with latent correlation 0.8 (the driver and five
redundant copies without an effect of their own), four blocks of eight distractor features
with correlation 0.7 and no effect, a pair of independent features whose product enters
the hazard, and independent noise features up to $p$. Features have Gaussian marginals.
Event times follow a Weibull proportional hazards model with shape 1.5, with the effect
size calibrated to a concordance of 0.70 before censoring, and censoring times are
independent and exponential, calibrated to the target fraction.

\paragraph{Evaluation.} Each dataset goes through the harness of the real cohorts, with
five outer folds and the ridge Cox model, and six selectors: \pkg{survfcbf} in permutation
mode and with tuned $q$ (chosen from $\{4, 7, 10\}$ by the untruncated Uno's C of the
downstream model over three inner folds, with 100 permutations), and the four baselines. The top-$K$ mode is evaluated on the same datasets at a matched panel size,
below. Table~\ref{S-tab:simulation}
reports two failure measures for the two self-sizing \pkg{survfcbf} modes, with their
Monte Carlo standard errors. The Application Note quotes the first: over all datasets,
the permutation mode returned an empty panel in every fold for \SuppSimEmptyPerm\% of
them.

\begin{table}[htbp]
\centering\footnotesize
\caption{Failures of the two self-sizing \pkg{survfcbf} modes in the simulation grid, as
percentages of the \SuppSimReps{} replicates of each configuration, with Monte Carlo
standard errors in brackets. Empty: no outer fold selected any feature. All failed: every
outer fold failed, because it selected fewer than the two features the downstream model
needs. Events: mean (Monte Carlo standard error). $^{\dagger}$Reference configuration.}
\label{S-tab:simulation}
\begin{tabular}{@{}rrr r rr rr@{}}
\toprule
 &  &  &  & \multicolumn{2}{c}{Permutation} & \multicolumn{2}{c}{Tuned $q$} \\
\cmidrule(lr){5-6}\cmidrule(l){7-8}
$n$ & Censored & $p$ & Events & Empty & All failed & Empty & All failed \\
\midrule
150 & 50\% & 2{,}000 & 76 (1.3) & 66.7 (8.6) & 100.0 (0.0) & 83.3 (6.8) & 93.3 (4.6) \\
300 & 50\% & 2{,}000 & 155 (1.5) & 70.0 (8.4) & 100.0 (0.0) & 53.3 (9.1) & 93.3 (4.6) \\
600$^{\dagger}$ & 50\% & 2{,}000 & 302 (2.8) & 70.0 (8.4) & 100.0 (0.0) & 10.0 (5.5) & 53.3 (9.1) \\
1{,}200 & 50\% & 2{,}000 & 600 (3.4) & 10.0 (5.5) & 36.7 (8.8) & 0.0 (0.0) & 0.0 (0.0) \\
2{,}400 & 50\% & 2{,}000 & 1{,}205 (5.7) & 0.0 (0.0) & 0.0 (0.0) & 0.0 (0.0) & 0.0 (0.0) \\
600 & 20\% & 2{,}000 & 479 (1.8) & 26.7 (8.1) & 86.7 (6.2) & 0.0 (0.0) & 3.3 (3.3) \\
600 & 70\% & 2{,}000 & 183 (2.6) & 83.3 (6.8) & 96.7 (3.3) & 26.7 (8.1) & 83.3 (6.8) \\
600 & 50\% & 200 & 300 (2.4) & 23.3 (7.7) & 93.3 (4.6) & 0.0 (0.0) & 6.7 (4.6) \\
600 & 50\% & 1{,}000 & 307 (2.2) & 46.7 (9.1) & 100.0 (0.0) & 3.3 (3.3) & 23.3 (7.7) \\
600 & 50\% & 5{,}000 & 303 (2.6) & 63.3 (8.8) & 96.7 (3.3) & 6.7 (4.6) & 53.3 (9.1) \\
\midrule
All &  &  &  & 46.0 (2.9) & 81.0 (2.3) & 18.3 (2.2) & 41.0 (2.8) \\
\bottomrule
\end{tabular}

\end{table}

\paragraph{Matched panel size.} \code{sim\_matched\_size.R} runs top-$K$ at $q = 2$, mRMR,
MIM and mRMR on the same two-bin target, and the three rankers on the same
\SuppSimDatasets{} datasets, folds and downstream model,
at $K = 3$, the number of driver clusters, and at $K = 10$. Each fold's selection is scored
against the cluster-level truth with \code{score\_selection\_clusterwise()}. A driver
cluster is recovered when any of its members is selected, the false discovery rate is the
share of distractor and noise features among those selected, copies are the members of a
driver cluster selected beyond the first, and exact recovery requires every driver cluster
and no distractor or noise feature, copies allowed (Table~\ref{S-tab:simmatched}). At
$K = 10$ at most three of the ten places can hold distinct driver clusters, so the false
discovery rate and exact recovery favour a selector that fills the panel with copies, and
only recovery and copies compare the selectors on equal terms there. Top-$K$ recovers the
most driver clusters at both sizes (\SuppSimTopkThree{} and \SuppSimTopkTen) and never
selects two members of one cluster. mRMR at its default binning also avoids copies, but
with a higher false discovery rate. mRMR on the same two-bin target comes within 0.02
(\SuppSimMrmrQtwoThree{} and \SuppSimMrmrQtwoTen) with almost no copies
(\SuppSimMrmrQtwoCopies{} per fold at $K = 10$), so most of the gain over the rankers comes
from the binning, not from pruning against averaging. MIM on that target recovers more
than at $q = 10$ but fills the panel with copies. The variance filter recovers almost
nothing, because the simulator decouples marginal variance from relevance, and its
downstream model fell back to the null model in \SuppSimMNullVar{} of \SuppSimMFolds{} folds
(\SuppSimMNullTopk{} for top-$K$). The simulator builds in the redundancy structure that the
pruning assumes, interchangeable clusters with one effect each, so these results show the
mechanism at work, not the advantage to expect on expression data, where
Section~\ref{S-sec:cardinality} finds no gain in redundancy over mRMR and CAR.

\begin{table}[htbp]
\centering\footnotesize
\caption{\pkg{survfcbf} top-$K$ ($q = 2$), mRMR, MIM and mRMR on the two-bin target, and
the rankers at a matched panel size $K$ on the \SuppSimDatasets{} simulated datasets, as
means with Monte Carlo standard errors in brackets. Recovered: share of the three driver
clusters with a member selected. FDR: share of the selected features that are distractors
or noise. Copies: members of a driver cluster selected beyond the first. Exact: share of
folds that recover every driver cluster with no distractor or noise feature. Uno's C:
truncated.}
\label{S-tab:simmatched}
\setlength{\tabcolsep}{4pt}
\begin{tabular}{@{}l r rrrrrr@{}}
\toprule
Selector & $K$ & Recovered & FDR & Copies & Exact & Uno's C & Stability \\
\midrule
\pkg{survfcbf} & 3 & 0.631 (0.018) & 0.369 (0.018) & 0.000 (0.000) & 0.353 (0.023) & 0.615 (0.003) & 0.408 (0.014) \\
mRMR & 3 & 0.315 (0.018) & 0.684 (0.018) & 0.000 (0.000) & 0.116 (0.017) & 0.565 (0.004) & 0.144 (0.013) \\
mRMR, $q = 2$ & 3 & 0.613 (0.018) & 0.386 (0.018) & 0.000 (0.000) & 0.339 (0.023) & 0.612 (0.003) & 0.384 (0.014) \\
MIM, $q = 2$ & 3 & 0.454 (0.012) & 0.260 (0.019) & 0.859 (0.032) & 0.054 (0.007) & 0.600 (0.003) & 0.438 (0.013) \\
Variance & 3 & 0.018 (0.003) & 0.975 (0.005) & 0.014 (0.006) & 0.000 (0.000) & 0.504 (0.002) & 0.237 (0.008) \\
CAR & 3 & 0.345 (0.014) & 0.453 (0.022) & 0.604 (0.035) & 0.045 (0.008) & 0.578 (0.003) & 0.387 (0.013) \\
MIM & 3 & 0.293 (0.015) & 0.608 (0.021) & 0.295 (0.024) & 0.045 (0.008) & 0.565 (0.003) & 0.181 (0.013) \\
\addlinespace
\pkg{survfcbf} & 10 & 0.776 (0.016) & 0.766 (0.005) & 0.000 (0.000) & 0.000 (0.000) & 0.615 (0.003) & 0.282 (0.006) \\
mRMR & 10 & 0.480 (0.019) & 0.854 (0.006) & 0.001 (0.001) & 0.000 (0.000) & 0.583 (0.004) & 0.098 (0.005) \\
mRMR, $q = 2$ & 10 & 0.762 (0.016) & 0.766 (0.005) & 0.042 (0.008) & 0.000 (0.000) & 0.614 (0.003) & 0.274 (0.006) \\
MIM, $q = 2$ & 10 & 0.718 (0.016) & 0.374 (0.018) & 4.106 (0.140) & 0.207 (0.020) & 0.624 (0.003) & 0.499 (0.013) \\
Variance & 10 & 0.049 (0.006) & 0.974 (0.004) & 0.092 (0.022) & 0.000 (0.000) & 0.508 (0.002) & 0.294 (0.005) \\
CAR & 10 & 0.548 (0.018) & 0.527 (0.020) & 3.069 (0.155) & 0.134 (0.017) & 0.603 (0.004) & 0.461 (0.013) \\
MIM & 10 & 0.470 (0.019) & 0.692 (0.018) & 1.654 (0.127) & 0.089 (0.016) & 0.588 (0.004) & 0.203 (0.013) \\
\bottomrule
\end{tabular}

\end{table}

\section{Validation against the reference implementation}\label{S-sec:validation}

The information-theoretic core of \pkg{survfcbf} was compared with the \pkg{FCBF} package
\citep{fcbfbioc}, distributed by Bioconductor until release 3.19, on classification data,
where both implement the same algorithm
(\code{fcbf\_crossvalidation.R}). Every check is a fixed
list of test cases drawn with a fixed seed. Panel A of Table~\ref{S-tab:validation}
compares the two packages. Symmetrical uncertainty agrees to $\SuppSuMaxDiff$ on random
pairs of factors. Selections, compared as the same features in the same order, agree on
every generated dataset without an exact SU tie. At an exact tie, when the SU between two
features equals the relevance of the less relevant one, the published $\ge$ rule
\citep{yu2003} removes that feature and the \pkg{FCBF} code, which tests $>$, keeps it.
Datasets with such a tie are therefore left out of that check. A second sweep of 300 new
datasets, drawn with another seed, allows duplicated features so that ties occur, and
every dataset in which the two selections differ contains an exact tie. The single-cell dataset scDengue
\citep{zanini2018}, shipped with \pkg{FCBF}, was discretised with that package's own
\code{discretize\_exprs()}, so that discretisation cannot cause a difference. Panel B checks
that the top-$K$ mode is an exact early stop, returning the first $K$ features of the
full run.

\begin{table}[htbp]
\centering\footnotesize
\caption{Validation of the information-theoretic core. Compared: datasets, thresholds or
values of $K$ compared. At a tie: mismatches at an exact SU tie. Skipped: generated
datasets left out of the check because they contain an exact tie.}
\label{S-tab:validation}
\begin{tabularx}{\linewidth}{@{}X rrrr@{}}
\toprule
Check & Compared & Mismatches & At a tie & Skipped \\
\midrule
\multicolumn{5}{@{}l}{\textit{A. \pkg{survfcbf} against the \pkg{FCBF} package}} \\
SU of random factor pairs ($n$ 20 to 120, 2 to 6 levels) & 200 & 0 & 0 &  \\
Selections, 400 generated datasets, those without an exact SU tie & 395 & 0 & 0 & 5 \\
Selections, 300 generated datasets with duplicated features allowed & 300 & 33 & 33 &  \\
Selections on scDengue (19{,}541 genes, 120 cells), five thresholds & 5 & 0 & 0 &  \\
\midrule
\multicolumn{5}{@{}l}{\textit{B. top-$K$ run against the first $K$ features of the full run}} \\
400 generated datasets, $K$ = 1 to 6, one of four thresholds each & 2{,}333 & 0 & 0 &  \\
With precomputed relevance and observation weights & 1{,}800 & 0 & 0 &  \\
Duplicated columns and tied relevance & 3{,}000 & 0 & 0 &  \\
scDengue, threshold 0.05, $K$ = 1 to 10 & 10 & 0 & 0 &  \\
\midrule
All top-$K$ checks & 7{,}143 & 0 &  &  \\
\bottomrule
\end{tabularx}

\end{table}

\section{Runtime}\label{S-sec:runtime}

Selection time was measured on one simulated dataset per value of $p$, with $n = 300$,
70\% censoring and the cluster structure of Section~\ref{S-sec:simulation}
(\code{runtime\_scaling.R}). Each time is the median of five elapsed wall-clock
measurements on one workstation, for \pkg{survfcbf} in top-$K$ mode ($K = 10$, $q = 2$)
and for the rankers cut at ten per cent of $p$ (Table~\ref{S-tab:runtime}). The scaling
exponent of \pkg{survfcbf}, \SuppRuntimeExponent, is the slope of a least-squares fit of
$\log$ time on $\log p$, over five values of $p$, so it is a coarse estimate: the cost is
close to linear in $p$. On real data the same selection took \SuppRealTimeMin{} to
\SuppRealTimeMax{} seconds per outer fold in the TCGA cohorts ($p$ from \SuppRealPMin{} to
\SuppRealPMax) and \SuppMetabricTime{} seconds in METABRIC, on the cluster.

\begin{table}[htbp]
\centering\footnotesize
\caption{Selection time in seconds by number of features, $n = 300$.}
\label{S-tab:runtime}
\begin{tabular}{@{}r rrrr@{}}
\toprule
$p$ & \pkg{survfcbf} top-$K$ & Variance & CAR & MIM \\
\midrule
1{,}000 & 0.44 & 0.01 & 0.46 & 0.20 \\
2{,}000 & 0.88 & 0.01 & 0.75 & 0.39 \\
5{,}000 & 2.12 & 0.03 & 1.79 & 1.09 \\
10{,}000 & 4.15 & 0.08 & 3.89 & 2.13 \\
20{,}000 & 8.18 & 0.15 & 10.44 & 10.94 \\
\bottomrule
\end{tabular}

\end{table}

\section{Software and computation}\label{S-sec:software}

All analyses use R \SuppLockR. The cross-validated benchmarks, the simulation grid and
the per-gene scores of the gene-set enrichment ran on the FinisTerrae III supercomputer
of CESGA, in its R 4.5.2 and Bioconductor 3.21 container with a separate package library.
In that library \pkg{mlr3}, \pkg{mlr3proba}, \pkg{mlr3extralearners}, \pkg{carSurv},
\pkg{praznik} and \pkg{stabm} have the versions of Table~\ref{S-tab:software}, and
\pkg{glmnet} is at 5.0. The other analyses, every table of this supplement and the
figure of the Application Note run on a workstation in the environment pinned
by the repository's \code{renv.lock} (R \SuppLockR, Bioconductor \SuppLockBioc), whose
main packages are listed in Table~\ref{S-tab:software}. The \pkg{mlr3proba} and \pkg{mlr3extralearners}
packages are not on CRAN, and the lockfile fixes them to their GitHub commits. Every
table, and the figure, regenerates from the result files tracked in the repository, without
access to the cluster or to the network. The repeated cross-validation of METABRIC, its
clinical reference and its censoring control ran on the workstation, with \pkg{glmnet}
4.1-10 instead of the 5.0 of the cluster. Its first split reproduces the selections of the
published METABRIC run exactly, and its Uno's C differs by at most
\SuppMetabricGlmnetDiff, the effect of the \pkg{glmnet} version on the ridge fit with about
two thousand features. The repeated cross-validation of TCGA, the matched panel sizes
with MIM and mRMR on the two-bin target, the random panels, the alternative target, the
centre analysis, the stability threshold, the groups of the pruning trace and the
simulation at matched size also ran on the cluster, in the same container and package
library, and \code{panel\_redundancy.R} saved the package versions of each job beside its
results.

\paragraph{Rerun with the Spearman tie-break.} The tie-break of
Section~\ref{S-sec:selectors} was added after the first runs. The \pkg{survfcbf} rows of
the TCGA benchmark (top-$K$ and permutation modes), of the matched panel size analysis and
of METABRIC, and the \pkg{survfcbf} scores of the gene-set enrichment, were then rerun
with it, and only those rows were replaced (\code{swap\_rerun\_rows.R}). This is valid
because a selector's row depends only on the data, the seed and its own code: the folds,
the per-fold seeds and the similarity matrix of the stability index do not depend on the
other selectors. The rerun included a control, top-$K$ with the old column-order rule,
which reproduced the stored top-$K$ rows exactly in every cohort and in METABRIC. Each
rerun job saved its \code{sessionInfo()} beside its results. The tuned-$q$ variant of the first runs on the real cohorts was not rerun, and its rows, still
in \code{benchmark\_realdata.csv}, are not reported. The simulation grid was not rerun
either. Instead, \code{grid\_tiebreak\_check.R} reran the permutation mode with both rules
at the same threshold on \SuppTieCheckReps{} replicates of each of the
\SuppTieCheckConfigs{} grid configurations (\SuppTieCheckDatasets{} datasets), on each
outer training set, at the default $q$ and at each $q$ of the tuned grid. All
\SuppTieCheckN{} selections were the same set in the same order. In none of them did two
features above the threshold tie on relevance (\SuppTieCheckMulti{} had at least two such
features), so at $q \geq 4$ the tie-break had nothing to decide in these data.

\begin{table}[htbp]
\centering\footnotesize
\caption{Main R packages, with the versions pinned in \code{renv.lock}.}
\label{S-tab:software}
\setlength{\tabcolsep}{4pt}
\begin{tabular}{@{}l l l l@{}}
\toprule
Package & Version & Source & Role \\
\midrule
\pkg{survival} & 3.8-3 & CRAN & Null Cox model, martingale residuals \\
\pkg{mlr3} & 1.7.1 & CRAN & Learner and measure framework \\
\pkg{mlr3proba} & 0.8.10 & GitHub mlr-org/mlr3proba@327d050 & Survival tasks, Uno's C, integrated Brier score \\
\pkg{mlr3extralearners} & 1.5.2 & GitHub mlr-org/mlr3extralearners@ad16c88 & Cross-validated penalised Cox learner \\
\pkg{glmnet} & 4.1-10 & CRAN & Ridge and lasso Cox models \\
\pkg{carSurv} & 1.0.0 & CRAN & CAR scores \\
\pkg{praznik} & 12.0.0 & CRAN & MIM and mRMR filters \\
\pkg{stabm} & 1.2.2 & CRAN & Selection stability \\
\pkg{fgsea} & 1.36.2 & Bioconductor & Gene-set enrichment \\
\pkg{msigdbr} & 26.1.0 & CRAN & MSigDB gene sets \\
\pkg{org.Hs.eg.db} & 3.22.0 & Bioconductor & Ensembl to symbol mapping \\
\pkg{FCBF} & 1.99.1 & GitHub lubianat/FCBF@51b9600 & Reference implementation for the validation \\
\bottomrule
\end{tabular}

\end{table}
